\input{preamble}

% OK, start here.
%
\begin{document}

\title{Introduction}


\maketitle

\phantomsection
\label{section-phantom}

\begin{verbatim}
Copyright (C) 2005 -- 2025 Johan de Jong
Permission is granted to copy, distribute and/or modify this
document under the terms of the GNU Free Documentation License,
Version 1.2 or any later version published by the Free Software
Foundation; with no Invariant Sections, no Front-Cover Texts,
and no Back-Cover Texts. A copy of the license is included in
the section entitled "GNU Free Documentation License".
\end{verbatim}

\tableofcontents


\section{Vue d'ensemble}
\label{section-overview}

\noindent
Outre le livre de Laumon et Moret-Bailly (voir \cite{LM-B}) et les travaux
(en cours) de Fulton et al., nous pensons qu'un manuel libre consacr\'e aux
champs alg\'ebriques et \`a la g\'eom\'etrie alg\'ebrique n\'ecessaire \`a leur
d\'efinition a sa place. Le projet Stacks tente d'y parvenir en construisant les
fondements \`a partir de l'alg\`ebre commutative, puis en passant par la
th\'eorie des sch\'emas et des espaces alg\'ebriques pour aboutir \`a des
fondements complets de la th\'eorie des champs alg\'ebriques.

\medskip\noindent
Nous pr\'evoyons que ce contenu sera lu en ligne, car l'une de ses
caract\'eristiques essentielles est la pr\'esence d'hyperliens donnant un
acc\`es rapide aux r\'ef\'erences internes r\'eparties sur de nombreuses pages.
Si vous utilisez dans votre navigateur une visionneuse PDF ou DVI int\'egr\'ee,
les liens entre fichiers devraient fonctionner.

\medskip\noindent
Ce projet est le fruit d'un travail collaboratif et nous vous encourageons
\`a y contribuer. Veuillez signaler par courriel toute coquille ou erreur
que vous relevez au cours de votre lecture, ainsi que toute suggestion,
tout contenu suppl\'ementaire ou tout exemple, \`a l'adresse
\href{mailto:stacks.project@gmail.com}{stacks.project@gmail.com}.
Vous pouvez t\'el\'echarger une archive tar contenant tous les fichiers sources,
l'extraire, ex\'ecuter la commande make et utiliser localement une visionneuse DVI ou PDF.
N'h\'esitez pas \`a modifier les fichiers LaTeX et \`a envoyer vos am\'eliorations
par courriel.


\section{Attribution}
\label{section-attribution}

\noindent
Vu l'ampleur de cet ouvrage, rendre justice, de façon à la fois exacte et
succincte, à tous les mathématiciens dont
les travaux ont conduit au d\'eveloppement de la th\'eorie que nous tentons
d'exposer ici est une tâche redoutable. Nous esp\'erons susciter \`a terme suffisamment d'int\'er\^et au
sein de la communaut\'e pour trouver des contributeurs dispos\'es \`a r\'ediger,
pour chacun des chapitres, des sections contenant des remarques historiques.

\medskip\noindent
Les personnes ayant contribu\'e \`a cet ouvrage sont mentionn\'ees sur la page
de titre de sa version sous forme de livre ainsi que
\href{https://stacks.math.columbia.edu/tex/CONTRIBUTORS}{sur cette page en ligne}.
Nous souhaitons ici signaler une s\'election de contributions majeures :
\begin{enumerate}
\item Jarod Alper a rédigé un chapitre consacré à la bibliographie relative
aux champs algébriques, voir
Guide bibliographique, section \ref{guide-section-short-introductions}.
\item Bhargav Bhatt a r\'edig\'e la premi\`ere version d'un chapitre
sur les morphismes \'etales de sch\'emas, voir
Morphismes \'etales, section \ref{etale-section-introduction}.
\item Bhargav Bhatt a r\'edig\'e la premi\`ere version de
Compl\'ements d'alg\`ebre, section \ref{more-algebra-section-formal-glueing}.
\item Kiran Kedlaya a fourni la premi\`ere r\'edaction de
Descente, section \ref{descent-section-descent-universally-injective}.
\item Les premi\`eres versions de
\begin{enumerate}
\item Alg\`ebre, section \ref{algebra-section-oka-families},
\item Objets injectifs, section \ref{injectives-section-baer}, et
\item le chapitre sur les corps, voir
Corps, section \ref{fields-section-introduction},
\end{enumerate}
proviennent du
\href{https://math.uchicago.edu/~amathew/cr.html}{The CRing Project},
avec l'aimable concours d'Akhil Mathew et al.
\item Alex Perry a r\'edig\'e le contenu consacr\'e aux modules projectifs
et aux modules de Mittag-Leffler, y compris la preuve de
Alg\`ebre, th\'eor\`eme \ref{algebra-theorem-ffdescent-projectivity}.
\item Alex Perry a r\'edig\'e le chapitre sur la th\'eorie des d\'eformations
selon l'approche de Schlessinger et Rim, voir
Th\'eorie formelle des d\'eformations, section \ref{formal-defos-section-introduction}.
\item Thibaut Pugin, Zachary Maddock et Min Lee ont pris des notes pour un
cours qui a servi de base \`a un chapitre sur la cohomologie \'etale
et \`a un chapitre sur la formule des traces. Voir
Cohomologie \'etale, section \ref{etale-cohomology-section-introduction}
et
La formule des traces, section \ref{trace-section-introduction}.
\item David Rydh a apport\'e de nombreux commentaires utiles, signal\'e plusieurs
erreurs, apport\'e une aide essentielle au contenu consacr\'e aux gerbes r\'esiduelles
et est \`a l'origine du contenu de
Compl\'ements sur les groupo\"ides dans les espaces, sections
\ref{spaces-more-groupoids-section-finite} et
\ref{spaces-more-groupoids-section-etale-localize}.
\item Burt Totaro a apport\'e les contributions suivantes : Exemples, sections
\ref{examples-section-non-descending-property-projective},
\ref{examples-section-non-effective-descent-projective},
ainsi que
Propri\'et\'es des champs algébriques, section \ref{stacks-properties-section-dimension}.
\item Le chapitre sur la cohomologie pro-\'etale, voir
Cohomologie pro-\'etale, section \ref{proetale-section-introduction},
est tir\'e d'un article de Bhargav Bhatt et Peter Scholze.
\item Bhargav Bhatt a apport\'e les contributions suivantes : Exemples, sections
\ref{examples-section-non-algebraic-hom-stack} et
\ref{examples-section-flat-not-colimit-flat-finitely-presented}.
\item Ofer Gabber a relev\'e des erreurs et apport\'e des corrections ; il a
\'egalement apport\'e les contributions suivantes :
Vari\'et\'es, lemme \ref{varieties-lemma-image-connected-component},
Espaces algébriques formels, lemme \ref{formal-spaces-lemma-completion-countably-indexed},
le contenu de
Compl\'ements sur les schémas en groupoïdes, section \ref{more-groupoids-section-ind-quasi-affine},
le r\'esultat principal de
Propri\'et\'es des espaces algébriques, section \ref{spaces-properties-section-fpqc},
ainsi que la preuve de
Compl\'ements sur la platitude, proposition
\ref{flat-proposition-finite-type-injective-into-flat-mod-m}.
\item J\'anos Koll\'ar a apport\'e les contributions suivantes :
Alg\`ebre, lemme \ref{algebra-lemma-hart-serre-loc-thm}, et
Cohomologie locale, proposition \ref{local-cohomology-proposition-kollar}.
\item Kiran Kedlaya a r\'edig\'e la premi\`ere version de
Compl\'ements d'alg\`ebre, section \ref{more-algebra-section-beauville-laszlo}.
\item Matthew Emerton, Toby Gee et Brandon Levin ont apport\'e
certains r\'esultats sur les \'epaississements, en particulier
Compl\'ements sur les morphismes de champs, lemmes
\ref{stacks-more-morphisms-lemma-reduced-diagonal},
\ref{stacks-more-morphisms-lemma-thickening-diagonals}, et
\ref{stacks-more-morphisms-lemma-thickening-properties}.
\item Lena Min Ji a r\'edig\'e la premi\`ere version de
Compl\'ements d'alg\`ebre, section \ref{more-algebra-section-principal-radical-ideals}.
\item Matthew Emerton et Toby Gee ont r\'edig\'e les premi\`eres versions de
G\'eom\'etrie des champs,
sections \ref{stacks-geometry-section-multiplicities} et
\ref{stacks-geometry-section-dimension-of-algebraic-stacks}.
\end{enumerate}



\begin{multicols}{2}[\section{Autres chapitres}]
\noindent
Préliminaires
\begin{enumerate}
\item \hyperref[introduction-section-phantom]{Introduction}
\item \hyperref[conventions-section-phantom]{Conventions}
\item \hyperref[sets-section-phantom]{Théorie des ensembles}
\item \hyperref[categories-section-phantom]{Catégories}
\item \hyperref[topology-section-phantom]{Topologie}
\item \hyperref[sheaves-section-phantom]{Faisceaux sur les espaces}
\item \hyperref[sites-section-phantom]{Sites et faisceaux}
\item \hyperref[stacks-section-phantom]{Champs}
\item \hyperref[fields-section-phantom]{Corps}
\item \hyperref[algebra-section-phantom]{Algèbre commutative}
\item \hyperref[brauer-section-phantom]{Groupes de Brauer}
\item \hyperref[homology-section-phantom]{Algèbre homologique}
\item \hyperref[derived-section-phantom]{Catégories dérivées}
\item \hyperref[simplicial-section-phantom]{Méthodes simpliciales}
\item \hyperref[more-algebra-section-phantom]{Compléments d'algèbre}
\item \hyperref[smoothing-section-phantom]{Lissage des morphismes d'anneaux}
\item \hyperref[modules-section-phantom]{Faisceaux de modules}
\item \hyperref[sites-modules-section-phantom]{Modules sur les sites}
\item \hyperref[injectives-section-phantom]{Objets injectifs}
\item \hyperref[cohomology-section-phantom]{Cohomologie des faisceaux}
\item \hyperref[sites-cohomology-section-phantom]{Cohomologie sur les sites}
\item \hyperref[dga-section-phantom]{Algèbre différentielle graduée}
\item \hyperref[dpa-section-phantom]{Algèbre à puissances divisées}
\item \hyperref[sdga-section-phantom]{Faisceaux différentiels gradués}
\item \hyperref[hypercovering-section-phantom]{Hyperrecouvrements}
\end{enumerate}
Schémas
\begin{enumerate}
\setcounter{enumi}{25}
\item \hyperref[schemes-section-phantom]{Schémas}
\item \hyperref[constructions-section-phantom]{Constructions de schémas}
\item \hyperref[properties-section-phantom]{Propriétés des schémas}
\item \hyperref[morphisms-section-phantom]{Morphismes de schémas}
\item \hyperref[coherent-section-phantom]{Cohomologie des schémas}
\item \hyperref[divisors-section-phantom]{Diviseurs}
\item \hyperref[limits-section-phantom]{Limites de schémas}
\item \hyperref[varieties-section-phantom]{Variétés}
\item \hyperref[topologies-section-phantom]{Topologies sur les schémas}
\item \hyperref[descent-section-phantom]{Descente}
\item \hyperref[perfect-section-phantom]{Catégories dérivées de schémas}
\item \hyperref[more-morphisms-section-phantom]{Compléments sur les morphismes}
\item \hyperref[flat-section-phantom]{Compléments sur la platitude}
\item \hyperref[groupoids-section-phantom]{Schémas en groupoïdes}
\item \hyperref[more-groupoids-section-phantom]{Compléments sur les schémas en groupoïdes}
\item \hyperref[etale-section-phantom]{Morphismes étales de schémas}
\end{enumerate}
Thèmes de théorie des schémas
\begin{enumerate}
\setcounter{enumi}{41}
\item \hyperref[chow-section-phantom]{Homologie de Chow}
\item \hyperref[intersection-section-phantom]{Théorie de l'intersection}
\item \hyperref[pic-section-phantom]{Schémas de Picard des courbes}
\item \hyperref[weil-section-phantom]{Théories cohomologiques de Weil}
\item \hyperref[adequate-section-phantom]{Modules adéquats}
\item \hyperref[dualizing-section-phantom]{Complexes dualisants}
\item \hyperref[duality-section-phantom]{Dualité pour les schémas}
\item \hyperref[discriminant-section-phantom]{Discriminants et différentes}
\item \hyperref[derham-section-phantom]{Cohomologie de de Rham}
\item \hyperref[local-cohomology-section-phantom]{Cohomologie locale}
\item \hyperref[algebraization-section-phantom]{Géométrie algébrique et formelle}
\item \hyperref[curves-section-phantom]{Courbes algébriques}
\item \hyperref[resolve-section-phantom]{Résolution des surfaces}
\item \hyperref[models-section-phantom]{Réduction semi-stable}
\item \hyperref[functors-section-phantom]{Foncteurs et morphismes}
\item \hyperref[equiv-section-phantom]{Catégories dérivées de variétés}
\item \hyperref[pione-section-phantom]{Groupes fondamentaux de schémas}
\item \hyperref[etale-cohomology-section-phantom]{Cohomologie étale}
\item \hyperref[crystalline-section-phantom]{Cohomologie cristalline}
\item \hyperref[proetale-section-phantom]{Cohomologie pro-étale}
\item \hyperref[relative-cycles-section-phantom]{Cycles relatifs}
\item \hyperref[more-etale-section-phantom]{Compléments de cohomologie étale}
\item \hyperref[trace-section-phantom]{La formule des traces}
\end{enumerate}
Espaces algébriques
\begin{enumerate}
\setcounter{enumi}{64}
\item \hyperref[spaces-section-phantom]{Espaces algébriques}
\item \hyperref[spaces-properties-section-phantom]{Propriétés des espaces algébriques}
\item \hyperref[spaces-morphisms-section-phantom]{Morphismes d'espaces algébriques}
\item \hyperref[decent-spaces-section-phantom]{Espaces algébriques décents}
\item \hyperref[spaces-cohomology-section-phantom]{Cohomologie des espaces algébriques}
\item \hyperref[spaces-limits-section-phantom]{Limites d'espaces algébriques}
\item \hyperref[spaces-divisors-section-phantom]{Diviseurs sur les espaces algébriques}
\item \hyperref[spaces-over-fields-section-phantom]{Espaces algébriques sur des corps}
\item \hyperref[spaces-topologies-section-phantom]{Topologies sur les espaces algébriques}
\item \hyperref[spaces-descent-section-phantom]{Descente et espaces algébriques}
\item \hyperref[spaces-perfect-section-phantom]{Catégories dérivées d'espaces}
\item \hyperref[spaces-more-morphisms-section-phantom]{Compléments sur les morphismes d'espaces}
\item \hyperref[spaces-flat-section-phantom]{Platitude sur les espaces algébriques}
\item \hyperref[spaces-groupoids-section-phantom]{Groupoïdes dans les espaces algébriques}
\item \hyperref[spaces-more-groupoids-section-phantom]{Compléments sur les groupoïdes dans les espaces}
\item \hyperref[bootstrap-section-phantom]{Amorçage}
\item \hyperref[spaces-pushouts-section-phantom]{Sommes amalgamées d'espaces algébriques}
\end{enumerate}
Thèmes de géométrie
\begin{enumerate}
\setcounter{enumi}{81}
\item \hyperref[spaces-chow-section-phantom]{Groupes de Chow des espaces}
\item \hyperref[groupoids-quotients-section-phantom]{Quotients de groupoïdes}
\item \hyperref[spaces-more-cohomology-section-phantom]{Compléments sur la cohomologie des espaces}
\item \hyperref[spaces-simplicial-section-phantom]{Espaces simpliciaux}
\item \hyperref[spaces-duality-section-phantom]{Dualité pour les espaces}
\item \hyperref[formal-spaces-section-phantom]{Espaces algébriques formels}
\item \hyperref[restricted-section-phantom]{Algébrisation des espaces formels}
\item \hyperref[spaces-resolve-section-phantom]{Résolution des surfaces revisitée}
\end{enumerate}
Théorie des déformations
\begin{enumerate}
\setcounter{enumi}{89}
\item \hyperref[formal-defos-section-phantom]{Théorie formelle des déformations}
\item \hyperref[defos-section-phantom]{Théorie des déformations}
\item \hyperref[cotangent-section-phantom]{Le complexe cotangent}
\item \hyperref[examples-defos-section-phantom]{Problèmes de déformation}
\end{enumerate}
Champs algébriques
\begin{enumerate}
\setcounter{enumi}{93}
\item \hyperref[algebraic-section-phantom]{Champs algébriques}
\item \hyperref[examples-stacks-section-phantom]{Exemples de champs}
\item \hyperref[stacks-sheaves-section-phantom]{Faisceaux sur les champs algébriques}
\item \hyperref[criteria-section-phantom]{Critères de représentabilité}
\item \hyperref[artin-section-phantom]{Axiomes d'Artin}
\item \hyperref[quot-section-phantom]{Espaces de Quot et de Hilbert}
\item \hyperref[stacks-properties-section-phantom]{Propriétés des champs algébriques}
\item \hyperref[stacks-morphisms-section-phantom]{Morphismes de champs algébriques}
\item \hyperref[stacks-limits-section-phantom]{Limites de champs algébriques}
\item \hyperref[stacks-cohomology-section-phantom]{Cohomologie des champs algébriques}
\item \hyperref[stacks-perfect-section-phantom]{Catégories dérivées de champs}
\item \hyperref[stacks-introduction-section-phantom]{Introduction aux champs algébriques}
\item \hyperref[stacks-more-morphisms-section-phantom]{Compléments sur les morphismes de champs}
\item \hyperref[stacks-geometry-section-phantom]{La géométrie des champs}
\end{enumerate}
Thèmes de théorie des espaces de modules
\begin{enumerate}
\setcounter{enumi}{107}
\item \hyperref[moduli-section-phantom]{Champs de modules}
\item \hyperref[moduli-curves-section-phantom]{Espaces de modules de courbes}
\end{enumerate}
Divers
\begin{enumerate}
\setcounter{enumi}{109}
\item \hyperref[examples-section-phantom]{Exemples}
\item \hyperref[exercises-section-phantom]{Exercices}
\item \hyperref[guide-section-phantom]{Guide bibliographique}
\item \hyperref[desirables-section-phantom]{Desiderata}
\item \hyperref[coding-section-phantom]{Style de programmation}
\item \hyperref[obsolete-section-phantom]{Obsolète}
\item \hyperref[fdl-section-phantom]{Licence de documentation libre GNU}
\item \hyperref[index-section-phantom]{Index généré automatiquement}
\end{enumerate}
\end{multicols}


\bibliography{my}
\bibliographystyle{amsalpha}

\end{document}
